\section{Theoretical Derivations}\label{app:proofs}

\subsection{Proof of Proposition~\ref{prop:joint_surrogate}}

\begin{proof}
The marginal refinement objective is
\begin{equation*}
    \KL\!\rbr{
        p_0\rbr{\dv}
        \,\|\,
         {  \int}
    K_{\theta,t}\rbr{\dv\mid\dv'}
    p_{\theta,t}\rbr{\dv'}
    \,d\dv'
    }
    =
    {  \int}
    p_0\rbr{\dv}
    \log
    \frac{
        p_0\rbr{\dv}
    }{
        {  \int}
        K_{\theta,t}\rbr{\dv\mid\dv'}
        p_{\theta,t}\rbr{\dv'}
        \,d\dv'
    }
    \,d\dv .
\end{equation*}
By Jensen's inequality and the concavity of $\log$, for every fixed $\dv$,
\begin{equation*}
    {  \int}
    q\rbr{\dv'|\dv}
    \log \frac{K_{\theta,t}\rbr{\dv\mid\dv'}
    p_{\theta,t}
    \rbr{\dv'}}{{q\rbr{\dv'|\dv}}}
    \,d\dv'
    \le
    \log
    {  \int}
    {q\rbr{\dv'|\dv}}\frac{K_{\theta,t}\rbr{\dv\mid\dv'}
    p_{\theta,t}
    \rbr{\dv'}}{{q\rbr{\dv'|\dv}}}
    \,d\dv' .
\end{equation*}
Multiplying by $p_0\rbr{\dv}$ and integrating over $\dv$ gives
\begin{equation*}
    \KL\!\rbr{
        p_0\rbr{\dv}
        \,\middle\|\,
        {  \int}
        K_{\theta,t}\rbr{\dv\mid\dv'}
        p_{\theta,t}\rbr{\dv'}
        \,d\dv'
    }
    \le
    \KL\!\rbr{
        p_0\rbr{\dv}q(\dv'|\dv)
        \,\middle\|\,
        K_{\theta,t}\rbr{\dv\mid\dv'}p_{\theta,t}\rbr{\dv'}
    }.
\end{equation*}
This proves the desired upper bound.
\end{proof}


\subsection{Derivation of Refinement Objective in Eq.~\eqref{eq:conditional_refinement_loss}}

We show that minimizing the joint surrogate reduces to conditional negative log-likelihood training. 
We have proven that the original matching objective in Eq.~\eqref{eq:ideal_refinement} is upper-bounded by the joint surrogate
\begin{equation*}
    \KL\!\rbr{
        p_0\rbr{\dv}p_{\theta,t}\rbr{\dv'}
        \,\middle\|\,
        K_{\theta,t}\rbr{\dv\mid\dv'}p_{\theta,t}\rbr{\dv'}
    }.
\end{equation*}
Expanding the KL divergence gives
\begin{align*}
    &
    \KL\!\rbr{
        p_0\rbr{\dv}p_{\theta,t}\rbr{\dv'|\dv}
        \,\middle\|\,
        K_{\theta,t}\rbr{\dv\mid\dv'}p_{\theta,t}\rbr{\dv'}
    }
    \nonumber\\
    &=
    {  \int}
    p_0\rbr{\dv}p_{\theta,t}\rbr{\dv'|\dv}
    \log
    \frac{
        p_0\rbr{\dv}p_{\theta,t}\rbr{\dv'|\dv}
    }{
        K_{\theta,t}\rbr{\dv\mid\dv'}p_{\theta,t}\rbr{\dv'}
    }
    \,d\dv\,d\dv'
    \nonumber\\
    &=
    {  \int}
    p_0\rbr{\dv}p_{\theta,t}\rbr{\dv'}
    \log \rbr{p_0\rbr{\dv}p_{\theta,t}\rbr{\dv'}}
    \,d\dv\,d\dv'
    -
    {  \int}
    p_0\rbr{\dv}p_{\theta,t}\rbr{\dv'|\dv}
    \log \rbr{K_{\theta,t}\rbr{\dv\mid\dv'}p_{\theta,t}\rbr{\dv'}}
    \,d\dv\,d\dv'
    \nonumber\\
    &
    \propto
    -
    \EE_{\dv\sim p_0,\,\dv'\sim p_{\theta,t}}
    \sbr{
        \log K_{\theta,t}\rbr{\dv\mid\dv'}
    }.
\end{align*}
Therefore, minimizing the joint surrogate over $K_{\theta,t}$ is equivalent to minimizing
\begin{equation*}
    \mathcal{L}_{t}\rbr{\theta}
    =
    -
    \EE_{\dv\sim p_0,\,\dv'\sim p_{\theta,t}}
    \sbr{
        \log K_{\theta,t}\rbr{\dv\mid\dv'}
    }.
\end{equation*}
This proves Eq.~\eqref{eq:conditional_refinement_loss}.

\subsection{Proof of Proposition~\ref{prop:bon_transition}}
\label{app:proof_bon_transition}

\begin{proof}
Given a draft $\dv'$, the self-refinement proposer draws $N$ candidates independently,
\begin{equation*}
    \dv_1,\ldots,\dv_N
    \sim
    K_{\theta}^{\mathrm{SR}}\rbr{\cdot\mid\dv'}.
\end{equation*}
Conditioned on the candidate set $\cbr{\dv_i}_{i=1}^{N}$, the Gumbel-Max selection rule selects candidate $j$ with probability
\begin{equation*}
    \PP\rbr{
        J=j
        \mid
        \cbr{\dv_i}_{i=1}^{N}
    }
    =
    \frac{
        \exp\rbr{\tau\scorer\rbr{\dv_j}}
    }{
        \sum_{\ell=1}^{N}
        \exp\rbr{\tau\scorer\rbr{\dv_\ell}}
    }.
    \label{eq:app_bon_selection_prob}
\end{equation*}
We first compute the contribution of the event that the first candidate equals a given sequence $\dv$ and is selected. 
Conditioning on $\dv_1=\dv$ and integrating over the remaining candidates gives
\begin{equation}
\begin{aligned}
    &
    \PP\rbr{
        \dv_1=\dv %,\ J=1
        \mid
        \dv'
    }
    \\
    &=
    { \int} \PP\rbr{\texttt{$\xb_1$ is the best among $\cbr{\dv_i}_{i=1}^N$}\mid \xb',\cbr{\dv_i}_{i=1}^N} K_{\theta,\mathrm{SR}}\rbr{\cbr{\dv_i}_{i=1}^N\mid\xb'} \,d\dv_{2:N}\\
    &=
    K_{\theta}^{\mathrm{SR}}\rbr{\dv\mid\dv'}
    {\int}
    \prod_{i=2}^{N}
    K_{\theta}^{\mathrm{SR}}\rbr{\dv_i\mid\dv'}
    \frac{
        \exp\rbr{\tau\scorer\rbr{\dv}}
    }{
        \exp\rbr{\tau\scorer\rbr{\dv}}
        +
        \sum_{j=2}^{N}
        \exp\rbr{\tau\scorer\rbr{\dv_j}}
    }
    \,d\dv_{2:N}
    \\
    &=
    K_{\theta}^{\mathrm{SR}}\rbr{\dv\mid\dv'}
    \EE_{\dv_{2:N}\sim K_{\theta}^{\mathrm{SR}}\rbr{\cdot\mid\dv'}}
    \sbr{
        \frac{
            \exp\rbr{\tau\scorer\rbr{\dv}}
        }{
            \exp\rbr{\tau\scorer\rbr{\dv}}
            +
            \sum_{j=2}^{N}
            \exp\rbr{\tau\scorer\rbr{\dv_j}}
        }
    }.
\end{aligned}
\label{eq:app_first_candidate_contribution}
\end{equation}
Therefore, Eq.~\eqref{eq:app_first_candidate_contribution} gives the induced best-of-$N$ refinement transition,
\begin{equation}
\begin{aligned}
    K_{\theta,\psi}^{\mathrm{BoN}}\rbr{\dv\mid\dv'}
    \propto K_{\theta}^{\mathrm{SR}}\rbr{\dv\mid\dv'}
    \EE_{\dv_{2:N}\sim K_{\theta}^{\mathrm{SR}}\rbr{\cdot\mid\dv'}}
    \sbr{
        \frac{
            \exp\rbr{\tau\scorer\rbr{\dv}}
        }{
            \exp\rbr{\tau\scorer\rbr{\dv}}
            +
            \sum_{j=2}^{N}
            \exp\rbr{\tau\scorer\rbr{\dv_j}}
        }
    }.
\end{aligned}
\end{equation}
This proves Proposition~\ref{prop:bon_transition}.
\end{proof}

\subsection{Derivation of Eq.~\eqref{eq:bon_loss}}
\label{app:derive_bon_loss}

We derive the tractable scorer objective by plugging the best-of-$N$ refinement transition into the conditional refinement loss. 
Let $\dv_1\sim p_0$ denote the positive target sequence, and let $\dv_2,\ldots,\dv_N$ be proposal candidates sampled from $K_{\theta}^{\mathrm{SR}}\rbr{\cdot\mid\dv'}$. 
From Eq.~\eqref{eq:app_first_candidate_contribution}, the first-candidate contribution to the best-of-$N$ transition is
\begin{equation*}
\begin{aligned}
    &
    K_{\theta}^{\mathrm{SR}}\rbr{\dv_1\mid\dv'}
    \EE_{\dv_{2:N}\sim K_{\theta}^{\mathrm{SR}}\rbr{\cdot\mid\dv'}}
    \sbr{
        \frac{
            \exp\rbr{\tau\scorer\rbr{\dv_1}}
        }{
            \exp\rbr{\tau\scorer\rbr{\dv_1}}
            +
            \sum_{j=2}^{N}
            \exp\rbr{\tau\scorer\rbr{\dv_j}}
        }
    } .
\end{aligned}
\label{eq:app_positive_transition_contribution}
\end{equation*}
Both $N$ and $K_{\theta}^{\mathrm{SR}}\rbr{\dv_1\mid\dv'}$ are independent of the scorer parameters during scorer training, since the proposer is treated with stop-gradient. 
Thus, the scorer-dependent part of the conditional refinement loss in Eq.~\eqref{eq:conditional_refinement_loss} is
\begin{equation}
\begin{aligned}
    -
    \EE_{\dv_1\sim p_0}\sbr{
    \log
    \EE_{\dv_{2:N}\sim K_{\theta}^{\mathrm{SR}}\rbr{\cdot\mid\dv'}}
    \sbr{
        \frac{
            \exp\rbr{\tau\scorer\rbr{\dv_1}}
        }{
            \exp\rbr{\tau\scorer\rbr{\dv_1}}
            +
            \sum_{j=2}^{N}
            \exp\rbr{\tau\scorer\rbr{\dv_j}}
        }
    }
    }.
\end{aligned}
\label{eq:app_bon_log_expectation}
\end{equation}
To obtain a tractable objective, we apply Jensen's inequality. 
Since $-\log(\cdot)$ is convex, for any positive random variable $R$,
\begin{equation*}
    -\log \EE\sbr{R}
    \le
    \EE\sbr{-\log R}.
    \label{eq:app_jensen_bon}
\end{equation*}
Using
\begin{equation*}
    R
    =
    \frac{
        \exp\rbr{\tau\scorer\rbr{\dv_1}}
    }{
        \exp\rbr{\tau\scorer\rbr{\dv_1}}
        +
        \sum_{j=2}^{N}
        \exp\rbr{\tau\scorer\rbr{\dv_j}}
    },
\end{equation*}
Eq.~\eqref{eq:app_bon_log_expectation} is upper-bounded by
\begin{equation*}
\begin{aligned}
    -
    \EE_{t,\,\dv_1\sim p_0,\,\cbr{\dv_j}_{j=2}^{N}\sim K_{\theta}^{\mathrm{SR}}\rbr{\cdot\mid\dv'}}\sbr{
    \log
    \frac{
        \exp\rbr{\tau\scorer\rbr{\dv_1}}
    }{
        \exp\rbr{\tau\scorer\rbr{\dv_1}}
        +
        \sum_{j=2}^{N}
        \exp\rbr{\tau\scorer\rbr{\dv_j}}
    }}.
\end{aligned}
\end{equation*}
Therefore, dropping constants independent of the scorer gives
\begin{equation*}
\begin{aligned}
  &\Lcal_{\mathrm{BoN}}\rbr{\psi}\\
  =&
  \EE_{t,\,\dv_1\sim p_0,\,\cbr{\dv_j}_{j=2}^{N}\sim K_{\theta}^{\mathrm{SR}}\rbr{\cdot\mid \dv'}}
  \sbr{
    -\tau\scorer\rbr{\dv_1}
    +
    \log
    \left(
        \exp\rbr{\tau\scorer\rbr{\dv_1}}
        +
        \sum_{j=2}^{N}
        \exp\rbr{\tau\scorer\rbr{\dv_j}}
    \right)
  }.
\end{aligned}
\end{equation*}
This is Eq.~\eqref{eq:bon_loss}.

\subsection{Properties and Justifications of Recursive Marginal Refinement}
\label{app:refinement_justification}

Although the learned sampler has the backward-chain form
\begin{equation*}
    p_\theta\rbr{\dv_{0:T}}
    =
    p_{\theta,T}\rbr{\dv_T}
    \prod_{t=1}^{T}
    K_{\theta,t}\rbr{\dv_{t-1}\mid\dv_t},
\end{equation*}
its learning procedure differs from conventional diffusion because the intermediate states are sampled from the model rather than generated by a predefined forward corruption path. 
We justify the recursive marginal refinement objective through two properties.


\begin{proposition}[Marginal consistency]
\label{prop:marginal_consistency_independent}
Let $\gamma_t\rbr{\dv,\dv'}=p_0\rbr{\dv}p_{\theta,t}\rbr{\dv'}$ be the product joint used in the refinement surrogate. 
If a refinement kernel $K$ satisfies
\begin{equation*}
    \KL\!\rbr{
        p_0\rbr{\dv}p_{\theta,t}\rbr{\dv'}
        \,\middle\|\,
        K\rbr{\dv\mid\dv'}p_{\theta,t}\rbr{\dv'}
    }
    =
    0,
\end{equation*}
then the induced refined marginal matches the target distribution:
\begin{equation*}
    {\int}
    K\rbr{\dv\mid\dv'}p_{\theta,t}\rbr{\dv'}\,d\dv'
    =
    p_0\rbr{\dv}.
\end{equation*}
More generally, if the joint surrogate is small, then the marginal refinement objective is also small:
\begin{equation*}
    \KL\!\rbr{
        p_0
        \,\middle\|\,
        Kp_{\theta,t}
    }
    \le
    \KL\!\rbr{
        p_0\rbr{\dv}p_{\theta,t}\rbr{\dv'}
        \,\middle\|\,
        K\rbr{\dv\mid\dv'}p_{\theta,t}\rbr{\dv'}
    }.
\end{equation*}
\end{proposition}

\begin{proof}
If the joint surrogate equals zero, then the two joint distributions are equal almost everywhere:
\begin{equation*}
    p_0\rbr{\dv}p_{\theta,t}\rbr{\dv'}
    =
    K\rbr{\dv\mid\dv'}p_{\theta,t}\rbr{\dv'}.
\end{equation*}
Integrating both sides over $\dv'$ gives
\begin{align*}
    {\int}
    p_0\rbr{\dv}p_{\theta,t}\rbr{\dv'}\,d\dv'
    &=
    {\int}
    K\rbr{\dv\mid\dv'}p_{\theta,t}\rbr{\dv'}\,d\dv' .
\end{align*}
Since $p_{\theta,t}$ is normalized, the left-hand side is $p_0\rbr{\dv}$. 
Therefore,
\begin{equation*}
    p_0\rbr{\dv}
    =
    {\int}
    K\rbr{\dv\mid\dv'}p_{\theta,t}\rbr{\dv'}\,d\dv',
\end{equation*}
which proves exact marginal consistency.
\end{proof}


\begin{proposition}[Ideal monotonic refinement]
\label{prop:ideal_monotonicity}
Let $\mathcal{K}$ be a refinement family that contains the identity kernel
\begin{equation*}
    I\rbr{\dv\mid\dv'}=\mathbf{1}\rbr{\dv=\dv'}.
\end{equation*}
Given the current draft marginal $p_{\theta,t}$, define the ideal refinement kernel as
\begin{equation*}
    K_t^\star
    =
    \arg\min_{K\in\mathcal{K}}
    \KL\!\rbr{
        p_0\rbr{\dv}
        \,\|\,
        {\int}
        K\rbr{\dv\mid\dv'}
        p_{\theta,t}\rbr{\dv'}
        \,d\dv'
    }.
\end{equation*}
Let the corresponding ideal refined marginal be
\begin{equation*}
    p_{\theta,t-1}^{\star}\rbr{\dv}
    =
    {\int}
    K_t^\star\rbr{\dv\mid\dv'}
    p_{\theta,t}\rbr{\dv'}
    \,d\dv'.
\end{equation*}
Then the ideal refinement step cannot increase the marginal KL to the target distribution:
\begin{equation*}
    \KL\!\rbr{
        p_0\rbr{\dv}
        \,\|\,
        p_{\theta,t-1}^{\star}\rbr{\dv}
    }
    \le
    \KL\!\rbr{
        p_0\rbr{\dv}
        \,\|\,
        p_{\theta,t}\rbr{\dv}
    }.
\end{equation*}
\end{proposition}

\begin{proof}
By definition of $K_t^\star$, we have
\begin{align*}
    \KL\!\rbr{
        p_0\rbr{\dv}
        \,\|\,
        p_{\theta,t-1}^{\star}\rbr{\dv}
    } \nonumber 
    =
    \min_{K\in\mathcal{K}}
    \KL\!\rbr{
        p_0\rbr{\dv}
        \,\|\,
        {\int}
        K\rbr{\dv\mid\dv'}
        p_{\theta,t}\rbr{\dv'}
        \,d\dv'
    }.
\end{align*}
Since $\mathcal{K}$ contains the identity kernel $I\rbr{\dv\mid\dv'}=\mathbf{1}\rbr{\dv=\dv'}$, the minimization over $\mathcal{K}$ includes the choice $K=I$. 
Therefore,
\begin{align*}
    \min_{K\in\mathcal{K}}
    \KL\!\rbr{
        p_0\rbr{\dv}
        \,\|\,
        {\int}
        K\rbr{\dv\mid\dv'}
        p_{\theta,t}\rbr{\dv'}
        \,d\dv'
    } \nonumber \\
   \le
    \KL\!\rbr{
        p_0\rbr{\dv}
        \,\|\,
        {\int}
        I\rbr{\dv\mid\dv'}
        p_{\theta,t}\rbr{\dv'}
        \,d\dv'
    }.
\end{align*}
The identity kernel preserves the current marginal:
\begin{equation*}
    {\int}
    I\rbr{\dv\mid\dv'}
    p_{\theta,t}\rbr{\dv'}
    \,d\dv'
    =
    p_{\theta,t}\rbr{\dv}.
\end{equation*}
Thus,
\begin{equation*}
    \KL\!\rbr{
        p_0\rbr{\dv}
        \,\|\,
        p_{\theta,t-1}^{\star}\rbr{\dv}
    }
    \le
    \KL\!\rbr{
        p_0\rbr{\dv}
        \,\|\,
        p_{\theta,t}\rbr{\dv}
    },
\end{equation*}
which proves the claim.
\end{proof}


\subsection{Interpretation of Empirical Best-of-$N$ Selection}
\label{app:log_ratio}

Let the proposal marginal induced by the current draft distribution and self-refinement proposer be
\begin{equation*}
    p_{\theta,t}^{\mathrm{SR}}\rbr{\dv}
    =
    {\int}
    K_{\theta}^{\mathrm{SR}}\rbr{\dv\mid\dv'}
    p_{\theta,t}\rbr{\dv'}
    \,d\dv'.
    \label{eq:app_sr_marginal}
\end{equation*}
The finite-candidate ranking objective contrasts positives $\dv_1\sim p_0$ with proposal samples $\dv_j\sim p_{\theta,t}^{\mathrm{SR}}$. 
At the population optimum, the score logit satisfies
\begin{equation*}
    \tau\scorer^\star\rbr{\dv}
    =
    \log p_0\rbr{\dv}
    -
    \log p_{\theta,t}^{\mathrm{SR}}\rbr{\dv}
    +
    C,
    \label{eq:app_marginal_log_ratio}
\end{equation*}
where $C$ is independent of $\dv$. 
Thus, the scorer estimates how much more likely a sequence is under the target distribution than under the proposal marginal.

This also motivates the empirical selection rule used in implementation. 
For candidates generated from a particular draft $\dv'$, ranking by
\begin{equation*}
    \log K_{\theta}^{\mathrm{SR}}\rbr{\dv\mid\dv'}
    +
    \tau\scorer\rbr{\dv}
\end{equation*}
combines the draft-conditioned proposal likelihood with a marginal correction toward the target distribution. 
This preserves the local preference of the proposer while using the scorer to favor globally more data-like candidates.